\documentclass[../C-analyse.tex]{subfiles}\begin{document}
\section{Backend}
Dette afsnit analyserer valg af backend-teknologi for Echo. Backenden er central for projektet, fordi den skal håndtere forretningslogik, rollebaseret adgangskontrol, opgavestatus, betalinger, Karma Koins, rewards, ratings og kommunikation med databasen. Valget af backend-teknologi påvirker derfor både systemets sikkerhed, udviklingshastighed, testbarhed og hvor let arkitekturen kan dokumenteres og videreudvikles.\\

Analysen sammenligner ASP.NET Core med C\#, Go og Node.js. De tre teknologier er relevante, fordi de alle kan bruges til moderne API-baserede websystemer, men de har forskellige styrker i forhold til struktur, performance, ecosystem og gruppens kompetencer.

\subsection{Krav}
Echo kræver en backend, der kan håndtere flere roller og sikre, at brugere kun kan udføre handlinger, de er autoriseret til. Det gælder eksempelvis oprettelse af rewards, tildeling af opgaver, overførsel af Karma Koins og betaling for udførte opgaver. Backenden skal derfor understøtte sikker authentication og authorization, samt gøre det let at placere adgangskontrol tæt på API'erne og forretningslogikken.\\

Backenden skal også kunne håndtere transaktioner og datakonsistens. Flere user stories kobler opgaver, saldo, transaktionshistorik og ratings sammen. Når en opgave markeres som udført, skal relaterede ændringer ske kontrolleret, så systemet ikke ender med forkert saldo, manglende transaktioner eller ratings uden gyldig opgavehistorik.\\

Derudover skal backenden kunne leve op til de ikke-funktionelle krav om API-responstid på under 500 ms under normale forhold, sikker sessionshåndtering med JWT, fejltolerance, dokumentation og testdækning. Teknologien skal derfor have stærk understøttelse af API-udvikling, databaseintegration, testing, logging og struktureret kode.\\

% Kriterier
For at analysere backend-teknologierne er der opstillet kriterier, som tager udgangspunkt i Echos funktionelle og ikke-funktionelle krav.

\begin{table}[H]
    \centering
    \renewcommand{\arraystretch}{1.25}
    \begin{tabularx}{\textwidth}{|X|c|X|}
        \hline
        \rowcolor{gray!30}
        \textbf{Kriterie} & \textbf{Vægt} & \textbf{Relevans for Echo} \\
        \hline
        Egnethed til projektet & 20\% & Teknologien skal passe til en API-backend med roller, opgaver, betaling, rewards og ratings. \\
        \hline
        Maintainability & 10\% & Backenden skal kunne opdeles i services/moduler med klar forretningslogik og testbar struktur. \\
        \hline
        Performance & 5\% & API-kald bør kunne besvares inden for 500 ms under normale forhold. \\
        \hline
        Kompetencer i gruppen & 10\% & Valget skal afspejle gruppens erfaring med C\#, Go og Node.js. \\
        \hline
        Ecosystem/tooling & 5\% & Teknologien bør have gode værktøjer til API'er, ORM, test, logging, dokumentation og deployment. \\
        \hline
        Fremtidig skalerbarhed & 5\% & Backenden skal kunne vokse med flere brugere, services og integrationer. \\
        \hline
        Authentication, authorization og sikkerhed & 15\% & Rollebaseret adgangskontrol, JWT og sikker håndtering af brugerdata er centrale krav. \\
        \hline
        Database- og transaktionsunderstøttelse & 15\% & Opgaver, saldo, betaling og Karma Koins kræver stabile databaseoperationer og transaktioner. \\
        \hline
        API-dokumentation og testbarhed & 10\% & Projektet kræver dokumenterede API'er og høj backend-testdækning. \\
        \hline
        Asynkrone processer og baggrundsjobs & 5\% & Email, notifikationer og eventuelle betalinger kan kræve asynkron behandling. \\
        \hline
        \rowcolor{gray!30}
        \textbf{Total} & \textbf{100\%} & \\
        \hline
    \end{tabularx}
    \caption{Vurderingskriterier for valg af backend-teknologi}
    \label{tab:backend-kriterier}
\end{table}

% Kræver biblatex (eller natbib) og kilder.bib
% Scorerne i parentes svarer til tabel \ref{tab:backend-scoring}

% Kræver biblatex (eller natbib) og kilder.bib
% Scorerne svarer til tabel \ref{tab:backend-scoring}

\subsection*{Egnethed til projektet (20\%)}
Alle tre teknologier kan bruges til at bygge web-API'er, men de passer forskelligt til Echo. ASP.NET Core leverer dependency injection, middleware, logging og konfiguration som en integreret del af frameworket~\cite{ms-di}. Det passer til et domæne, hvor opgaver, betalinger, saldo og rewards naturligt kan opdeles i services med tydelige ansvarsområder. Node.js er en event-drevet runtime, der er velegnet til API'er med mange I/O-operationer~\cite{node-event-loop}. Det passer til Echos integrationer som emails og notifikationer, og hvis gruppen vælger Angular i frontend, giver en TypeScript-backend ét fælles sprog på tværs af systemet. Go er designet til netværksservices, og standardbiblioteket understøtter siden Go 1.22 routing på HTTP-metoder og path-parametre~\cite{go-routing}. Go overlader dog flere arkitekturvalg til gruppen og ville være mere oplagt i en distribueret arkitektur med mange små services end i et domænetungt system som Echo.

\medskip\noindent\textbf{Score:} ASP.NET Core: 9 \quad Go: 7 \quad Node.js: 8

\subsection*{Maintainability (10\%)}
C\# og Go er begge statisk typede, så mange fejl fanges ved kompilering frem for i drift. ASP.NET Core skiller sig ud ved den indbyggede dependency injection, der holder komponenter løst koblede, og ved en forudsigelig udgivelsesplan. Den aktuelle LTS-version, .NET 10, understøttes frem til november 2028~\cite{dotnet-support}. Go er bevidst simpelt, \texttt{gofmt} sikrer ensartet kodestil, og Go 1-kompatibilitetsløftet betyder, at eksisterende kode fortsat kan kompileres med nye versioner~\cite{go1compat}. Til gengæld findes der færre etablerede konventioner for at strukturere større forretningsapplikationer. Node.js står svagest. JavaScript er dynamisk typet, og Node.js' indbyggede TypeScript-understøttelse fjerner kun typerne uden at typetjekke koden~\cite{node-typescript}. Typesikkerhed kræver derfor et separat build-trin. Express har desuden ingen faste konventioner, så strukturen afhænger af gruppens disciplin.

\medskip\noindent\textbf{Score:} ASP.NET Core: 9 \quad Go: 7 \quad Node.js: 6

\subsection*{Performance (5\%)}
I TechEmpower Framework Benchmarks Round 23 ligger Go-frameworks blandt de hurtigste, og ASP.NET Core placerer sig ligeledes højt, mens udbredte Node.js-frameworks generelt ligger lavere~\cite{techempower-r23}. Go kompileres til native kode, og goroutines er så lette, at mange samtidige forespørgsler kan håndteres med lavt hukommelsesforbrug~\cite{effective-go}. Node.js' event-loop er effektivt til I/O, men da JavaScript-koden kører i én tråd, blokerer CPU-tungt arbejde alle andre forespørgsler~\cite{node-event-loop}. Benchmarkene er syntetiske og måler frameworkenes øvre grænse. Da Echos arbejde overvejende er databasebundet, vil alle tre levere tilstrækkelig performance i praksis. Det afspejles i kriteriets lave vægt.

\medskip\noindent\textbf{Score:} ASP.NET Core: 9 \quad Go: 10 \quad Node.js: 7

\subsection*{Kompetencer i gruppen (10\%)}
Gruppen har relativt høj erfaring med C\#, hvilket reducerer projektets risiko og frigør tid til arkitektur, test og dokumentation. Kompetencerne i Node.js er lavere. Erfaring med JavaScript eller TypeScript fra frontend giver et vist udgangspunkt, men backend-specifikke mønstre som middleware og ORM-valg kræver oplæring. Go er det mindst kendte sprog, og enkelte i gruppen har slet ikke arbejdet med det. I et tidsbegrænset bachelorprojekt er den nødvendige oplæring en betydelig risiko.

\medskip\noindent\textbf{Score:} ASP.NET Core: 8 \quad Go: 4 \quad Node.js: 5

\subsection*{Ecosystem/tooling (5\%)}
Node.js er den mest anvendte webteknologi i Stack Overflows udviklerundersøgelse fra 2025, og ASP.NET Core bruges af omkring en femtedel af respondenterne~\cite{so-survey-2025}. Begge har derfor store communities og biblioteker til stort set ethvert behov via henholdsvis npm og NuGet. .NET har desuden stærke IDE'er som Visual Studio og Rider. Go leveres med det stærkeste indbyggede tooling af de tre: \texttt{go build}, \texttt{go test}, \texttt{go vet}, \texttt{gofmt} og modulsystemet, og \texttt{govulncheck} kan scanne for kendte sårbarheder i afhængigheder~\cite{govulncheck}. Til gengæld er Go's biblioteksudvalg til forretningsnære funktioner som ORM og brugerstyring mindre og mere fragmenteret.

\medskip\noindent\textbf{Score:} ASP.NET Core: 9 \quad Go: 8 \quad Node.js: 9

\subsection*{Fremtidig skalerbarhed (5\%)}
Alle tre kan køre som stateless services i containere og skaleres horisontalt. Go har en fordel, fordi sproget kompileres til små, selvstændige binærfiler med hurtig opstart og lavt hukommelsesforbrug. ASP.NET Core har typisk højere hukommelsesforbrug og længere opstartstid, selvom Native AOT-kompilering kan mindske forskellen~\cite{ms-native-aot}. Node.js håndterer mange samtidige forbindelser effektivt, men én proces udnytter kun én CPU-kerne til JavaScript-kode. Der skal derfor skaleres med flere processer, f.eks. via cluster-modulet eller flere containere~\cite{node-cluster}.

\medskip\noindent\textbf{Score:} ASP.NET Core: 8 \quad Go: 9 \quad Node.js: 8

\subsection*{Authentication, authorization og sikkerhed (15\%)}
Her er forskellen mellem teknologierne størst. ASP.NET Core Identity håndterer brugere, passwords og login, og JWT-baseret authentication understøttes direkte~\cite{ms-identity}. Autorisation kan defineres både rollebaseret og policy-baseret~\cite{ms-authz-roles, ms-authz-policies}, hvilket passer til Echos behov for at styre, hvilke roller der må oprette opgaver, godkende betalinger og ændre saldo. Node.js har modne biblioteker som Passport.js og NestJS' guards, men løsningen sammensættes af tredjepartspakker. Det øger eksponeringen for supply chain-angreb. Go's standardbibliotek indeholder kryptografiske primitiver, men ingen samlet løsning til brugerstyring og roller~\cite{go-std}. Gruppen skal derfor selv skrive mere sikkerhedskritisk kode, hvilket øger risikoen for fejl i et system med betalinger.

\medskip\noindent\textbf{Score:} ASP.NET Core: 9 \quad Go: 6 \quad Node.js: 7

\subsection*{Database- og transaktionsunderstøttelse (15\%)}
Transaktionssikkerhed er central for Echo, hvor ændring af saldo og registrering i transaktionshistorikken skal ske atomisk. Alle tre teknologier understøtter databasetransaktioner, men integrationen er forskellig. Entity Framework Core giver migrations, relationer og LINQ-forespørgsler, og et kald til \texttt{SaveChanges} udføres som standard i én transaktion~\cite{ef-transactions}. Eksplicitte transaktioner kan bruges, når flere operationer skal lykkes eller fejle samlet. I Node.js understøtter ORM'er som Prisma og TypeORM relationer, migrations og transaktioner~\cite{prisma-transactions}, men ingen af dem er lige så tæt integreret med frameworket som EF Core. I Go understøtter \texttt{database/sql} transaktioner med \texttt{BeginTx}, \texttt{Commit} og \texttt{Rollback}~\cite{go-transactions}. Understøttelsen er solid, men lavniveau, og der findes ingen officiel ORM eller migrationsløsning. I både Node.js og Go falder valget og konfigurationen af værktøjer derfor på gruppen.

\medskip\noindent\textbf{Score:} ASP.NET Core: 9 \quad Go: 7 \quad Node.js: 7

\subsection*{API-dokumentation og testbarhed (10\%)}
Fra .NET 9 kan ASP.NET Core generere OpenAPI-dokumentation uden tredjepartsbiblioteker~\cite{ms-openapi}. Integrationstest understøttes via \texttt{WebApplicationFactory}, som starter API'et i hukommelsen~\cite{ms-integration-tests}, og dependency injection gør det let at bruge mocks i unit tests. Go har høj testbarhed, da \texttt{testing}- og \texttt{httptest}-pakkerne er en del af standardbiblioteket~\cite{go-testing, go-httptest}, men OpenAPI-dokumentation kræver tredjepartsværktøjer som swaggo eller oapi-codegen. I Node.js afhænger det af frameworkvalget. NestJS kan generere OpenAPI-dokumentation automatisk ud fra controllere og DTO'er~\cite{nestjs-openapi}, og Node.js har en stabil indbygget test runner~\cite{node-test}. Med Express kræver dokumentation derimod mere manuelt arbejde, og Node.js' score kunne derfor hæves til 8, hvis NestJS blev valgt.

\medskip\noindent\textbf{Score:} ASP.NET Core: 9 \quad Go: 7 \quad Node.js: 7

\subsection*{Asynkrone processer og baggrundsjobs (5\%)}
Echo har behov for baggrundsopgaver som afsendelse af emails, notifikationer og periodisk beregning af rewards. C\# har \texttt{async}/\texttt{await} indbygget, og ASP.NET Core understøtter baggrundsopgaver via \texttt{IHostedService} og \texttt{BackgroundService}~\cite{ms-hosted-services}. Til planlagte og persistente jobs findes modne biblioteker som Hangfire og Quartz.NET. I Go er goroutines og channels en del af selve sproget, så baggrundsopgaver kan startes med få linjer kode og afbrydes kontrolleret via \texttt{context}~\cite{effective-go}. Node.js er stærk til asynkron I/O, men langvarige eller CPU-tunge jobs blokerer event-loopet og kræver worker threads~\cite{node-worker-threads} eller en ekstern jobkø som BullMQ. For alle tre kræver jobs, der skal overleve en genstart, en persistent kø.

\medskip\noindent\textbf{Score:} ASP.NET Core: 9 \quad Go: 9 \quad Node.js: 8

\subsection{Scoring}
Scoringen er baseret på en skala fra 1 til 10, hvor 10 er bedst. Den vægtede score er beregnet ud fra kriteriets vægt.

\begin{table}[H]
    \centering
    \renewcommand{\arraystretch}{1.2}
    \begin{tabularx}{\textwidth}{|X|c|c|c|c|}
        \hline
        \rowcolor{gray!30}
        \textbf{Kriterie} & \textbf{Vægt} & \textbf{ASP.NET Core} & \textbf{Go} & \textbf{Node.js} \\
        \hline
        Egnethed til projektet & 20\% & 9 & 7 & 8 \\
        \hline
        Maintainability & 10\% & 9 & 7 & 6 \\
        \hline
        Performance & 5\% & 9 & 10 & 7 \\
        \hline
        Kompetencer i gruppen & 10\% & 8 & 4 & 5 \\
        \hline
        Ecosystem/tooling & 5\% & 9 & 8 & 9 \\
        \hline
        Fremtidig skalerbarhed & 5\% & 8 & 9 & 8 \\
        \hline
        Authentication, authorization og sikkerhed & 15\% & 9 & 6 & 7 \\
        \hline
        Database- og transaktionsunderstøttelse & 15\% & 9 & 7 & 7 \\
        \hline
        API-dokumentation og testbarhed & 10\% & 9 & 7 & 7 \\
        \hline
        Asynkrone processer og baggrundsjobs & 5\% & 9 & 9 & 8 \\
        \hline
        \rowcolor{gray!30}
        \textbf{Vægtet score} & \textbf{100\%} & \textbf{8,85} & \textbf{6,95} & \textbf{7,10} \\
        \hline
    \end{tabularx}
    \caption{Scoring af backend-kandidater}
    \label{tab:backend-scoring}
\end{table}

\subsection*{Beslutning}
På baggrund af analysen vurderes ASP.NET Core med C\# som det mest hensigtsmæssige backend-valg for Echo. Teknologien passer godt til en API-baseret platform med rollebaseret adgangskontrol, transaktioner, dokumentation og testbar forretningslogik. Samtidig har gruppen de stærkeste kompetencer i C\#, hvilket gør valget mere realistisk inden for projektets tidsramme.\\

ASP.NET Core giver også et godt grundlag for rapporten, fordi frameworkets struktur gør det let at beskrive services, controllers, dependency injection, databaseadgang og sikkerhed. Valget understøtter derfor både MVP-udviklingen og projektets akademiske fokus på proces og tekniske beslutninger.

\end{document}
